\scanpage{4}
classe \inserted{extérieure} d'\deleted{isom.} \unsure{exactement}] et la condition est

\diagram{\begin{tikzcd}[ampersand replacement=\&,row sep=tiny,column sep=tiny]
\& \deleted{\unsure{x_i}}\arrow[dl,dash,"\deleted{G_k}"']\arrow[dr,dash,"\deleted{G_j}"]\&\\
{}\arrow[rr,dash,"\deleted{G_i}"']\&\&{}
\end{tikzcd}}
\diagram{\begin{tikzcd}[ampersand replacement=\&,row sep=small,column sep=small]
G_k\arrow[dr,"{\tau_{ik}(x_j)}"']\&\&G_j\arrow[ll,"{\tau_{kj}(x_i)}"']\\
\&G_i\arrow[ur,"{\tau_{ji}(x_k)}"']\&
\end{tikzcd}}
\[
\begin{aligned}
&\tau_{ji}(x_k)\tau_{ik}(x_j)\tau_{kj}(x_i)\\
&\hspace{4em}=\mathrm{id}_{G_j}
\end{aligned}
\]
(condition qui est bien invariante par permutations circulaires)

L'\uline{existence} d'un tel $x=(x_i)_{i\in I}$ équivaut à la condition que les isom. \uline{extérieurs} $\widetilde\tau_{ij}=\widetilde\tau_{ij}(x_k)$ (\deleted{$i,j,k\in I$} $i,j,k\in I$ distincts) définis par les $E_k$ ($k\in I$) entre les $G_i$ forment un syst. transitif d'isom de liens.

La condition envisagée sur $x$ ne dépend que de son orbite modulo $\prod G_i$, \deleted{Soit}\inserted{qui} est aussi son orbite sous l'un quelconque des sous-gr. $\prod_{j\in I-\{i\}}G_j$.

\margin{\(\begin{aligned}
g_i\gamma_i x&=(g_j,g_k)x\\
g_i(\gamma_i,\gamma_j,\gamma_k)x&=(\gamma_j,\gamma_k)g_i\gamma_i x=(\gamma_j g_j,\gamma_k g_k)x
\end{aligned}\)}
Soit $\Gamma$ une telle orbite, qui est donc un torseur sous l'un quelconque des \inserted{trois} groupes $\prod_{j\in I-\{i\}}G_j$. On a ainsi \inserted{une} \deleted{\unsure{certaine|certain|contenu}} partie
\[
\Gamma\subset\prod_{i\in I} E_i=P
\]
Celle-ci satisfait à \inserted{la} condition de \textcircled{I}.

\[
\begin{aligned}
a&=(a_i,a_j,a_k)&x_i&=a_i g_k^{-1}& &\text{détermine }g_k\text{ en fonction de }a\text{ et }x_i\\
x&=(x_i\ \ x_j\ \ \underset{?}{x_k})&x_j&=g_k a_j g_i^{-1}& &\text{détermine }g_i\text{ en fonction de }a\text{ et }x_i,x_j\\
&&x_k&=g_i a_k&&\text{détermine }x_k
\end{aligned}
\]
% archive sheet number: 3
